\documentclass[10pt,aspectratio=169]{beamer}
% Preámbulo modular para presentaciones Beamer (Modelación Estadística)
% Diseño y paleta institucional del Tecnológico de Monterrey

\documentclass[aspectratio=169,xcolor={dvipsnames,table}]{beamer}

\usepackage[utf8]{inputenc}
\usepackage[spanish,mexico]{babel}
\usepackage{amsmath,amssymb,amsfonts,mathtools}
\usepackage{graphicx}
\usepackage{listings}
\usepackage{booktabs}
\usepackage{tikz}
\usetikzlibrary{matrix,arrows,decorations.pathmorphing,shapes.geometric,positioning}

% Paleta Institucional Tecnológico de Monterrey
\definecolor{TecAzul}{HTML}{0039A6}       % Azul TEC: color primario institucional
\definecolor{TecAzulOscuro}{HTML}{1A2E51} % Azul marino oscuro
\definecolor{TecRojo}{HTML}{EC2661}       % Rojo de acento (paleta miTec)
\definecolor{TecGris}{HTML}{646464}       % Gris neutro para comentarios y subtítulos
\definecolor{TecFondoBloque}{HTML}{F4F6F9}% Fondo suave para bloques

% Configuración de Tema Beamer
\usetheme{Boadilla}
\usecolortheme[named=TecAzulOscuro]{structure}
\setbeamercolor{alerted text}{fg=TecRojo}
\setbeamercolor{palette primary}{bg=TecAzulOscuro,fg=white}
\setbeamercolor{palette secondary}{bg=TecAzul,fg=white}
\setbeamercolor{palette tertiary}{bg=TecAzulOscuro!80!black,fg=white}
\setbeamercolor{title}{fg=TecAzulOscuro,bg=TecFondoBloque}
\setbeamercolor{frametitle}{fg=TecAzulOscuro,bg=TecFondoBloque}

% Bloques personalizados elegantes
\setbeamercolor{block title}{bg=TecAzulOscuro,fg=white}
\setbeamercolor{block body}{bg=TecFondoBloque,fg=black}
\setbeamercolor{block title alerted}{bg=TecRojo,fg=white}
\setbeamercolor{block body alerted}{bg=TecRojo!8,fg=black}
\setbeamercolor{block title example}{bg=TecAzul,fg=white}
\setbeamercolor{block body example}{bg=TecAzul!8,fg=black}

% Redondear bordes y añadir sombra sutil
\setbeamertemplate{blocks}[rounded][shadow=true]
\setbeamertemplate{navigation symbols}{}

% Pie de página limpio con número de diapositiva
\setbeamertemplate{footline}{
  \leavevmode%
  \hbox{%
  \begin{beamercolorbox}[wd=.7\paperwidth,ht=2.25ex,dp=1ex,left]{author in head/foot}%
    \hspace*{2ex}\insertshortauthor~~(\insertshortinstitute) \hfill \insertshorttitle
  \end{beamercolorbox}%
  \begin{beamercolorbox}[wd=.3\paperwidth,ht=2.25ex,dp=1ex,right]{date in head/foot}%
    \insertshortdate{}\hspace*{2em}
    \insertframenumber{} / \inserttotalframenumber\hspace*{2ex}
  \end{beamercolorbox}}%
  \vskip0pt%
}

% Configuración de Listings de Python para visualización pedagógica de código
\lstset{
    language=Python,
    basicstyle=\ttfamily\footnotesize,
    keywordstyle=\color{TecAzulOscuro}\bfseries,
    stringstyle=\color{TecRojo},
    commentstyle=\color{TecGris}\itshape,
    showstringspaces=false,
    breaklines=true,
    frame=lines,
    rulecolor=\color{TecAzulOscuro!30},
    backgroundcolor=\color{TecFondoBloque},
    aboveskip=0.5em,
    belowskip=0.5em,
    literate={á}{{\'a}}1 {é}{{\'e}}1 {í}{{\'i}}1 {ó}{{\'o}}1 {ú}{{\'u}}1
             {Á}{{\'A}}1 {É}{{\'E}}1 {Í}{{\'I}}1 {Ó}{{\'O}}1 {Ú}{{\'U}}1
             {ñ}{{\~n}}1 {Ñ}{{\~N}}1 {¿}{{?`}}1 {¡}{{!`}}1
}

% Metadatos institucionales por defecto
\title[Modelación Estadística]{Modelación Estadística para Ciencia de Datos e Ingeniería}
\author[J. Castillo Colmenares]{Juliho David Castillo Colmenares}
\institute[Tec de Monterrey]{Tecnológico de Monterrey \\ Escuela de Ingeniería y Ciencias}
\date{\today}
% Comandos y macros estadísticas compartidas para las presentaciones Beamer (Metropolis)

% Conjuntos numéricos básicos
\newcommand{\R}{\mathbb{R}}
\newcommand{\N}{\mathbb{N}}
\newcommand{\Q}{\mathbb{Q}}
\newcommand{\Z}{\mathbb{Z}}
\newcommand{\set}[1]{\left\{ #1 \right\}}
\newcommand{\abs}[1]{\left|#1\right|}
\newcommand{\norm}[1]{\left\| #1 \right\|}

% Comandos estadísticos y probabilísticos del curso
\newcommand{\Var}{\operatorname{Var}}
\newcommand{\cov}{\operatorname{Cov}}
\newcommand{\corr}{\rho}
\newcommand{\comb}[2]{\begin{pmatrix} #1 \\ #2 \end{pmatrix}}
\newcommand{\s}{\sigma}
\newcommand{\particion}{\mathfrak{P}}
\newcommand{\card}[1]{\#\left( #1 \right)}
\newcommand{\seq}[3]{\set{#1}_{#2}^{#3}}
\newcommand{\E}{\mathbb{E}}
\newcommand{\Prob}{\mathbb{P}}

% Entornos pedagógicos y cajas en español académico natural
\newenvironment{discusion}{%
  \begin{alertblock}{Pregunta de Discusión}%
}{%
  \end{alertblock}%
}

\newenvironment{reflexion}{%
  \begin{alertblock}{Reflexión para el Aula}%
}{%
  \end{alertblock}%
}

\newenvironment{paradoja}{%
  \begin{alertblock}{Paradoja de Exploración}%
}{%
  \end{alertblock}%
}

\newenvironment{motivacion}{%
  \begin{exampleblock}{Motivación: Ciencia de Datos en la Práctica}%
}{%
  \end{exampleblock}%
}

\newenvironment{retoclase}{%
  \begin{block}{Reto Rápido en Equipo}%
}{%
  \end{block}%
}

\title[Pruebas de varianzas]{Sección 07.08: Pruebas relacionadas con varianzas}\subtitle{Dispersión y homoscedasticidad}
\author[J. Castillo Colmenares]{Juliho Castillo Colmenares}\institute{Tecnológico de Monterrey}\date{}
\begin{document}
\begin{frame}[plain]\titlepage\end{frame}
\begin{frame}{Hoja de ruta}\small\begin{enumerate}\item Una varianza.\item Comparación de dos varianzas.\item Supuestos e interpretación.\end{enumerate}\end{frame}
\begin{frame}{Fuente y alcance}\small La sección presenta pruebas de dispersión para evaluar uniformidad, riesgo y supuestos de homoscedasticidad.\begin{block}{Pregunta guía}¿La variabilidad observada coincide con la variabilidad hipotética?\end{block}\end{frame}
\begin{frame}{Varianza poblacional}\small Para una muestra normal de tamaño $n$, la varianza muestral $S^2$ resume la dispersión alrededor de $\bar X$.\end{frame}
\begin{frame}{Hipótesis de una varianza}\small\[H_0:\sigma^2=\sigma_0^2\qquad\text{frente a una alternativa sobre }\sigma^2.\]\end{frame}
\begin{frame}{Estadística chi-cuadrada}\small\[\chi^2=\frac{(n-1)S^2}{\sigma_0^2}.\]\begin{block}{Distribución} Bajo $H_0$, $\chi^2\sim\chi^2_{n-1}$.\end{block}\end{frame}
\begin{frame}{Regiones de rechazo}\small La alternativa determina si se usa una cola derecha, izquierda o ambas colas de la distribución chi-cuadrada.\end{frame}
\begin{frame}{Dos varianzas}\small Para dos muestras normales independientes, se contrasta $H_0:\sigma_1^2=\sigma_2^2$ mediante el cociente de varianzas muestrales.\end{frame}
\begin{frame}{Estadística $F$}\small\[F=\frac{S_1^2}{S_2^2}.\]\begin{block}{Distribución} Bajo igualdad, $F\sim F_{n_1-1,n_2-1}$.\end{block}\end{frame}
\begin{frame}{Orden del cociente}\small Colocar la varianza mayor en el numerador facilita una prueba de cola derecha; los grados de libertad deben seguir el orden elegido.\end{frame}
\begin{frame}{Procedimiento I}\small\begin{enumerate}\item Revisar normalidad e independencia.\item Formular hipótesis.\item Fijar $\alpha$.\end{enumerate}\end{frame}
\begin{frame}{Procedimiento II}\small\begin{enumerate}\setcounter{enumi}{3}\item Calcular $S^2$ o $F$.\item Obtener cuantil o valor $p$.\item Concluir en contexto.\end{enumerate}\end{frame}
\begin{frame}{Supuestos}\small Las pruebas $\chi^2$ y $F$ son sensibles a desviaciones de normalidad. La independencia y el diseño muestral son esenciales.\end{frame}
\begin{frame}{Aplicación: una varianza}\small Suponga una referencia $\sigma_0^2$ y una muestra con $n$ y $S^2$. Calcule $\chi^2=(n-1)S^2/\sigma_0^2$.\end{frame}
\begin{frame}{Aplicación: decisión}\small Compare el valor con los cuantiles de $\chi^2_{n-1}$ según la alternativa.\end{frame}
\begin{frame}{Aplicación: dos varianzas}\small Dos procesos producen $S_1^2$ y $S_2^2$. La razón $F$ evalúa si sus dispersiones pueden considerarse iguales.\end{frame}
\begin{frame}{Aplicación: interpretación}\small Rechazar igualdad de varianzas señala dispersión distinta; no identifica por sí sola la causa ni el mecanismo del proceso.\end{frame}
\begin{frame}{Homoscedasticidad}\small En modelos paramétricos, la igualdad de varianzas sostiene errores comparables entre grupos. La prueba es una herramienta diagnóstica, no un sustituto del análisis gráfico.\end{frame}
\begin{frame}{Errores frecuentes}\small\begin{itemize}\item Omitir la normalidad.\item Invertir grados de libertad al reordenar $F$.\item Interpretar dispersión como diferencia de medias.\end{itemize}\end{frame}
\begin{frame}{Plantilla de reporte}\small\begin{block}{Reporte} ``Se obtuvo $\chi^2/F=\cdots$, con $gl=\cdots$ y $p=\cdots$. Al nivel $\alpha=\cdots$, [rechazamos/no rechazamos] la igualdad planteada.''\end{block}\end{frame}
\begin{frame}{Síntesis}\small\begin{itemize}\item Una varianza usa $\chi^2$.\item Dos varianzas usan $F$.\item Los supuestos controlan la validez de la inferencia.\end{itemize}\end{frame}
\begin{frame}{Conexión con el capítulo}\small Estas pruebas fundamentan la evaluación de homoscedasticidad en comparaciones de medias y modelos de regresión.\end{frame}
\end{document}
